\chapter{Contexte de l'entreprise et étude de l'existant}

\section{Présentation de l'entreprise}

CHO Group est une entreprise tunisienne opérant dans le secteur agroalimentaire, spécialisée dans la production et l'exportation de l'huile d'olive. Fondée en 1996, l'entreprise s'est progressivement développée pour devenir un acteur important sur les marchés internationaux.

Le groupe est présent dans plusieurs pays et se distingue par un modèle intégré couvrant l'ensemble de la chaîne de valeur, allant de la production agricole jusqu'à la commercialisation des produits finis.

\begin{figure}[H]
\centering

\includegraphics[width=0.35\textwidth]{cho_group_logo.png}
\caption{Logo de CHO Group}
\label{fig:cho_logo_ch1}
\end{figure}

\begin{figure}[H]
\centering
\includegraphics[width=0.90\textwidth]{Filiales.png}
\caption{Présence internationale de CHO Group}
\label{fig:cho_filiales_ch1}
\end{figure}

Cette dimension internationale implique une forte exposition à la concurrence ainsi qu'à la variabilité des prix, ce qui rend le suivi des marchés particulièrement stratégique pour l'entreprise.

\section{Activités de l'entreprise}

Les activités principales de CHO Group couvrent l'ensemble du processus de production et de distribution :

\begin{itemize}[leftmargin=1.2cm]
	\item production d'huile d'olive ;
	\item raffinage et transformation ;
	\item conditionnement ;
	\item exportation vers plusieurs marchés internationaux.
\end{itemize}

\begin{figure}[H]
\centering
\includegraphics[width=0.90\textwidth]{un_group_integre.png}
\caption{Chaîne de valeur et processus intégré de CHO Group}
\label{fig:cho_integre_ch1}
\end{figure}

Ces activités impliquent la gestion de marchés multiples et nécessitent un suivi rigoureux des prix afin de maintenir la compétitivité et optimiser les décisions commerciales.

\section{Étude de l'existant}

Avant la mise en place de ce projet, le processus de suivi des prix au sein de l'entreprise reposait principalement sur des méthodes manuelles.

Dans un premier temps, les équipes se déplacent vers les points de vente (marchés, magasins, grandes surfaces) afin de collecter les prix des produits. En parallèle, certaines données peuvent être récupérées à partir de sources en ligne selon les besoins.

\begin{figure}[H]
\centering
\includegraphics[width=0.92\textwidth]{example_source_of_data.png}
\caption{Exemple de source de données utilisée pour la collecte des prix}
\label{fig:source_donnees_ch1}
\end{figure}

Dans un second temps, les données collectées sont saisies manuellement dans des fichiers Excel.

\begin{figure}[H]
\centering
\includegraphics[width=0.92\textwidth]{Excel_image_example.png}
\caption{Exemple de fichier Excel utilisé dans le processus actuel}
\label{fig:excel_existant_ch1}
\end{figure}

Enfin, les opérations de tri, de nettoyage et d'analyse sont réalisées manuellement, ce qui rend le processus long et dépendant des interventions humaines.

Le processus global peut être résumé comme suit :

\begin{itemize}[leftmargin=1.2cm]
	\item collecte terrain et/ou web ;
	\item saisie manuelle des données ;
	\item traitement et analyse manuelle.
\end{itemize}

\begin{figure}[H]
\centering
\includegraphics[width=0.92\textwidth]{old_company_solution.png}
\caption{Schéma du système actuel : collecte, Excel et analyse manuelle}
\label{fig:process_actuel_ch1}
\end{figure}

\section{Critique de l'existant}

L'analyse du système actuel met en évidence plusieurs limites importantes :

\begin{itemize}[leftmargin=1.2cm]
	\item un processus lent, dû à la collecte manuelle et aux traitements non automatisés ;
	\item un risque élevé d'erreurs lors de la saisie des données ;
	\item une difficulté d'accès rapide à une information fiable et consolidée ;
	\item une absence de vision globale et en temps réel des prix ;
	\item une forte dépendance au travail humain.
\end{itemize}

Ces contraintes impactent directement la qualité et la rapidité de la prise de décision au sein de l'entreprise.

\section{Objectifs du projet}

Afin de répondre aux limites identifiées, ce projet a pour objectif de mettre en place une solution de Business Intelligence permettant :

\begin{itemize}[leftmargin=1.2cm]
	\item d'automatiser la collecte des données à partir de différentes sources ;
	\item de centraliser les données dans une base unique ;
	\item d'améliorer la qualité et la fiabilité des données ;
	\item de produire des indicateurs clés (KPI) pour le pilotage ;
	\item de fournir des tableaux de bord interactifs pour l'analyse.
\end{itemize}

Ce projet se concentre sur une approche descriptive et analytique, constituant une base solide pour de futures évolutions vers des analyses plus avancées.

\section{Solution proposée}

La solution proposée consiste en une plateforme de Business Intelligence permettant de transformer un processus manuel en un système automatisé et structuré.

Elle repose sur plusieurs composants principaux :

\begin{itemize}[leftmargin=1.2cm]
	\item des sources de données (web et terrain) ;
	\item un processus ETL pour la collecte et le traitement ;
	\item une base de données centralisée ;
	\item une interface de visualisation sous forme de tableaux de bord.
\end{itemize}

\begin{figure}[H]
\centering
\includegraphics[width=0.92\textwidth]{our_diagram_solution.png}
\caption{Schéma de la solution BI : sources, ETL, base de données et dashboard}
\label{fig:solution_bi_ch1}
\end{figure}

Cette solution permet de réduire l'intervention manuelle, d'améliorer la fiabilité des données et de fournir une vision claire et exploitable des prix, facilitant ainsi la prise de décision.

\section{Conclusion}

Ce chapitre a permis de présenter le contexte général de l'entreprise ainsi que le fonctionnement du système actuel. L'analyse de l'existant a mis en évidence plusieurs limites justifiant la mise en place d'une solution de Business Intelligence adaptée.

Le chapitre suivant sera consacré à l'analyse des besoins et à la conception de la solution proposée.

\clearpage
